% !TEX root = chapter-standalone.tex

\section[Free-forgetful adjunction]{Examples of ``Free-Forgetful'' adjunctions}
\label{sec:free-forgetful-adjunction-graph-example}

Another ``type'' of \SY{adjunction} that appears frequently can be called a ``Free-Forgetful'' \SY{adjunction}.
Such \SY{adjunctions} are composed of a ``free functor'' and a ``forgetful functor''.
These terms are informal, but the idea is as follows.

A free \SY{functor}~$\CatC \fto \CatD$ typically takes an object~$\Obja$ of~\CatC and ``freely'' adds some structure to it.
``Free'' means that only those structures and conditions are added that are absolutely necessary to make~$\Obja$ an object of~\CatD, and otherwise the \SY{functor} does not impose any further constraints or relations.

Conversely, a ``forgetful functor'' usually starts from an object~$\Objb$ in~\CatD which has some structure, and ``forgets'' some of this structure, which results in us being able to view~$\Objb$ as an object in~\CatC.

\begin{example}
    Any real \SY{vector space} is built from an underlying set, together with extra structure given by operations (vector addition and scalar multiplication).
    There is a forgetful \SY{functor}
    \begin{equation}\label{eq:forgetful-functor-from-VectR}
        \VectR \fto \Set
    \end{equation}
    which maps any \SY{vector space} to its underlying set of vectors.
    On the other hand, there is a ``free'' functor
    \begin{equation}\label{eq:free-functor-to-VectR}
        \Set \fto \VectR.
    \end{equation}
    Given a set~\setA, we can build the ``free real \SY{vector space} generated by~\setA''.
    To do this, we think of the elements of~\setA as basis vectors, and we build a \SY{vector space} by taking formal finite~\reals-linear combinations of them.
\end{example}

In the following we consider an example in detail where we ``freely'' generate a category from a directed graph.

\begin{example}
    Let~\Graph be the category of (directed, multi-) graphs and~\Category the category of (small) categories.

    There is a \SY{functor}~$\funa \colon \Graph \fto \Category$ which turns any graph~$G = \tup{V,E, s,t}$ into a category whose objects are the vertices~$V$ and whose morphisms are finite directed paths between vertices.
    This is called the \emph{free category generated by the graph~$G$} (\cref{sec:catsfromgraphs}).

    There is also a \SY{functor}~$\funb \colon \Category \fto \Graph$ which turns a category~C into a graph where the set of vertices is~$\Obof{C}$ and there is a directed edge between vertices for every morphism in C between the corresponding vertices.

    Let's first describe this \SY{adjunction} via \cref{def:adj-iso}.
    The natural isomorphism
    \begin{equation}\label{eq:cat-graph-adjunction-transpose-iso}
        \stylenat{\tau}\colon \HomSet{\Category}{\funa(-)}{-} \nto \HomSet{\Graph}{-}{\funb(-)}
    \end{equation}
    is the one whose component at~$\tup{G,C}$ is the isomorphism
    \begin{equation}\label{eq:cat-graph-adjunction-transpose-iso-components}
        \tau_{G,C} \colon \HomSet{\Category}{\funa(G)}{C} \nto \HomSet{\Graph}{G}{\funb(C)}
    \end{equation}
    which assigns to any \SY{functor}~$\funa\colon \funa(G) \fto C$ the morphism of graphs~$G \colon \funb(C)$ given by restricting~$\funa$ to~$G$ and only keeping track of its action on vertices and edges (in other words, we ignore its compositional properties and think of it just as a graph morphism).

    Now let's consider this \SY{adjunction} from the perspective of \cref{def:adj-counit}.
    The component at~$G$ of the counit is the morphism of graphs
    \begin{equation}\label{eq:cat-graph-adjunction-unit}
        \equivunit_D \colon G \to \funb(\funa(G))
    \end{equation}
    which includes~$G$ into the graph~$\funb(\funa(G))$.
    The latter has an edge from the source to the target of every finite path in~$G$.
    The paths of length zero are what corresponded to \SY{identity morphisms} in~$\funa(G)$, and the paths of length one constitute a copy of~$G$ inside~$\funb(\funa(G))$.

    What does the counit look like?
    Its component at~C is a functor
    \begin{equation}\label{eq:cat-graph-adjunction-co-unit}
        \equivcounit_{C} \colon \funa(\funb(C)) \fto C.
    \end{equation}

    The category~$\funa(\funb(C))$ is larger than~C: starting with C, the graph~$\funb(C)$ will contain edges for all the morphisms in~C, but it will forget their compositional interlinking.
    In particular, for example, it will forget which loops denote \SY{identity morphisms} (in other words, which morphisms act neutrally) and, more generally, it will forget when different compositions of morphism give the same result.
    In~$\funa(\funb(C))$, then, morphism compositions that might have given the same result in~C will now be distinct.
    The \SY{functor}~$\equivcounit_{C}$ in a sense ``remembers'' those relations that were true in C and it ``implements'' them by ``projecting''~$\funa(\funb(C))$ back to~C.
\end{example}
